\documentclass[14pt,a4paper,oneside]{extbook}
\usepackage[utf8]{vietnam}
\usepackage{amsmath,amssymb,mathrsfs,fancyhdr,enumerate,multirow,makecell}
\usepackage{tkz-euclide,tikz-3dplot,tikz,tkz-tab}
\usetikzlibrary{arrows.meta,calc,intersections,angles,quotes,backgrounds,shapes.geometric,patterns,through,shadows,positioning}
\usepackage{yhmath}
\usepackage[top=0.7cm, bottom=0.7cm, left=1.0cm, right=1.0cm, headheight=16pt, headsep=3mm, footskip=5mm]{geometry}
\usepackage[hidelinks,unicode]{hyperref}
\usepackage{qrcode}
\usepackage[many]{tcolorbox}
\usepackage{fontawesome}

\pagestyle{fancy}
\fancyhf{}
\fancyhead[L]{\footnotesize\bfseries\color{teal!85!black}\faBook~SỔ TAY KIẾN THỨC TOÁN 5}
\fancyhead[R]{\footnotesize\bfseries\color{blue!85!black}Lớp Toán Thầy Lĩnh -- loptoanthaylinh.vn}
\fancyfoot[C]{\footnotesize\bfseries\color{gray!80!black}-- Sổ tay kiến thức: Bài 9. Luyện tập chung --}
\renewcommand{\headrulewidth}{0.4pt}

\begin{document}

% --- BANNER TIÊU ĐỀ BÀI HỌC ---
\noindent
\begin{tcolorbox}[
  enhanced,
  colback=teal!8!white,
  colframe=teal!80!black,
  arc=3mm,
  boxrule=1.1pt,
  width=\linewidth,
  halign=center,
  top=2.2mm, bottom=2.2mm,
  left=3mm, right=3mm,
  drop fuzzy shadow
]
  {\large\bfseries\color{teal!95!black}\faBookmark\ BÀI 09. LUYỆN TẬP CHUNG (CHỦ ĐỀ 1)}\\[2pt]
  {\scriptsize\color{black!75}\itshape Tổng hợp kỹ năng tính giá trị biểu thức, tính nhanh và giải toán có lời văn trên 1 trang A4}
\end{tcolorbox}

\vspace{3.2mm}
% --- KHỐI 1: TÍNH GIÁ TRỊ BIỂU THỨC CHỨA PHÂN SỐ ---
\noindent
\begin{minipage}[t]{0.54\textwidth}
  \begin{tcolorbox}[
    enhanced,
    title={\footnotesize\bfseries\color{white}\faCalculator\ 1. Quy tắc tính giá trị biểu thức},
    colback=blue!2!white,
    colframe=blue!75!black,
    arc=2.5mm, boxrule=0.8pt, top=1.8mm, bottom=1.8mm, left=2.2mm, right=2.2mm
  ]
    \footnotesize
    $\bullet$ \textbf{Thứ tự ưu tiên:}\\
    \quad Trong ngoặc $\rightarrow$ Nhân, chia $\rightarrow$ Cộng, trừ.\\
    $\bullet$ Biểu thức chỉ có cộng, trừ (hoặc nhân, chia): tính từ \textbf{trái sang phải}.\\
    $\bullet$ Nếu có hỗn số: \textbf{chuyển đổi hết ra phân số} trước khi bắt đầu tính toán.
  \end{tcolorbox}
\end{minipage}%
\hfill
\begin{minipage}[t]{0.44\textwidth}
  \begin{tcolorbox}[
    enhanced,
    title={\footnotesize\bfseries\color{white}\faBolt\ 2. Kỹ thuật tính thuận tiện},
    colback=teal!2!white,
    colframe=teal!75!black,
    arc=2.5mm, boxrule=0.8pt, top=1.8mm, bottom=1.8mm, left=2mm, right=2mm
  ]
    \footnotesize
    $\bullet$ Đặt \textbf{thừa số chung}:\\
    \quad $A \times B + A \times C = A \times (B + C)$.\\
    $\bullet$ Ghép cặp tạo số nguyên:\\
    \quad $\left(\dfrac{2}{7} + \dfrac{5}{7}\right) = 1$; $\left(\dfrac{3}{8} + \dfrac{5}{8}\right) = 1$.\\
    $\bullet$ Trừ một tổng: $A - (B + C) = A - B - C$.
  \end{tcolorbox}
\end{minipage}

\vspace{3mm}

% --- KHỐI 2: HAI BÀI TOÁN CƠ BẢN VỀ PHÂN SỐ ---
\noindent
\begin{tcolorbox}[
  enhanced,
  title={\footnotesize\bfseries\color{white}\faCheckSquareO\ 3. Hai bài toán phân số kinh điển cần phân biệt rạch ròi},
  colback=purple!2!white,
  colframe=purple!70!black,
  arc=2.5mm, boxrule=0.8pt, top=1.8mm, bottom=1.8mm, left=2.5mm, right=2.5mm
]
  \footnotesize
  \begin{tabular}{p{0.46\textwidth}|p{0.48\textwidth}}
    \textbf{Bài toán 1: Tìm phân số của một số} & \textbf{Bài toán 2: Tìm một số khi biết phân số của nó} \\
    $\triangleright$ Muốn tìm $\dfrac{m}{n}$ của số $A$, ta lấy: & $\triangleright$ Muốn tìm số đó khi biết $\dfrac{m}{n}$ của nó bằng $B$, ta lấy: \\
    \centering\textbf{\color{purple!90!black}$A \times \dfrac{m}{n}$} & \centering\textbf{\color{blue!90!black}$B : \dfrac{m}{n}$} \\[2pt]
    \textit{VD:} Tìm $\dfrac{2}{3}$ của $18\text{ kg} \rightarrow 18 \times \dfrac{2}{3} = 12\text{ kg}$. & \textit{VD:} Tìm quãng đường biết $\dfrac{2}{3}$ quãng đường là $12\text{ km} \rightarrow 12 : \dfrac{2}{3} = 18\text{ km}$.
  \end{tabular}
\end{tcolorbox}

\vspace{3mm}

% --- KHỐI 3: SƠ ĐỒ TƯ DUY ---
\noindent
\begin{tcolorbox}[
  enhanced,
  title={\footnotesize\bfseries\color{white}\faSitemap\ Sơ đồ tư duy: Giải toán có lời văn về phân số},
  colback=yellow!2!white,
  colframe=teal!85!black,
  arc=3mm, boxrule=0.9pt, top=1.8mm, bottom=1.8mm, left=1.5mm, right=1.5mm,
  halign=center
]
  \begin{tikzpicture}[
    scale=0.92,
    every node/.style={transform shape},
    s1/.style={rectangle, rounded corners=4pt, fill=blue!15, draw=blue!70!black, text=blue!90!black, font=\bfseries\scriptsize, inner sep=4pt, align=center},
    s2/.style={rectangle, rounded corners=4pt, fill=teal!15, draw=teal!70!black, text=teal!90!black, font=\bfseries\scriptsize, inner sep=4pt, align=center},
    s3/.style={rectangle, rounded corners=4pt, fill=orange!15, draw=orange!70!black, text=orange!90!black, font=\bfseries\scriptsize, inner sep=4pt, align=center},
    s4/.style={rectangle, rounded corners=4pt, fill=green!15, draw=green!70!black, text=green!90!black, font=\bfseries\scriptsize, inner sep=4pt, align=center}
  ]
    \node[s1] (n1) at (-6.3, 0) {1. Đọc kĩ đề, tóm tắt\\(Vẽ sơ đồ đoạn thẳng)};
    \node[s2] (n2) at (-2.1, 0) {2. Xác định dạng toán\\(Tìm phần hay tìm tổng)};
    \node[s3] (n3) at (2.1, 0) {3. Thực hiện phép tính\\(Nhân hay Chia)};
    \node[s4] (n4) at (6.3, 0) {4. Kiểm tra \& Đáp số\\(Kèm đơn vị đo)};

    \draw[teal!85!black, thick, -{Stealth[length=2.5mm]}] (n1) -- (n2);
    \draw[teal!85!black, thick, -{Stealth[length=2.5mm]}] (n2) -- (n3);
    \draw[teal!85!black, thick, -{Stealth[length=2.5mm]}] (n3) -- (n4);
  \end{tikzpicture}
\end{tcolorbox}

\vspace{3mm}

% --- KHỐI 4: KỸ NĂNG VẼ SƠ ĐỒ & BẪY TRẮC NGHIỆM ---
\noindent
\begin{minipage}[t]{0.54\textwidth}
  \begin{tcolorbox}[
    enhanced,
    title={\footnotesize\bfseries\color{white}\faPencil\ Vũ khí vẽ sơ đồ đoạn thẳng},
    colback=green!3!white,
    colframe=green!60!black,
    arc=2.5mm, boxrule=0.8pt, top=1.8mm, bottom=1.8mm, left=2.2mm, right=2.2mm
  ]
    \footnotesize
    Sơ đồ đoạn thẳng giúp trực quan hóa bài toán:\\
    $\bullet$ Mẫu số chỉ \textbf{tổng số phần bằng nhau}.\\
    $\bullet$ Tử số chỉ \textbf{số phần chiếm giữ}.\\
    $\triangleright$ Tìm giá trị của $1$ phần $\rightarrow$ nhân với số phần cần tìm.
  \end{tcolorbox}
\end{minipage}%
\hfill
\begin{minipage}[t]{0.44\textwidth}
  \begin{tcolorbox}[
    enhanced,
    title={\footnotesize\bfseries\color{white}\faExclamationTriangle\ Mẹo tránh bẫy trắc nghiệm},
    colback=red!2!white,
    colframe=red!70!black,
    arc=2.5mm, boxrule=0.8pt, top=1.8mm, bottom=1.8mm, left=2.2mm, right=2.2mm
  ]
    \footnotesize
    $\bullet$ \textbf{Cực kỳ dễ nhầm}: Tìm $\dfrac{2}{3}$ của $A$ là làm phép \textbf{nhân} ($A \times \dfrac{2}{3}$), còn biết $\dfrac{2}{3}$ của một số là $B$ thì phải làm phép \textbf{chia} ($B : \dfrac{2}{3}$).
  \end{tcolorbox}
\end{minipage}

\end{document}
